% Exemplo: alinhamento da lista de referências no biblatex-abnt. A NBR 6023 pede referências
% alinhadas à esquerda; a opção justify do biblatex-abnt justifica a lista (fora da norma).
% Página 1: justify; página 2: padrão do estilo (\raggedright, sem hifenização); página 3: alinhado
% à esquerda com hifenização (ragged2e), URL na fonte do texto e quebra de URL (xurl).
% A mancha estreita só exagera as diferenças. Compilar: latexmk -pdf alinhamento-referencias.tex
% Testado com: abntex2, biblatex-abnt, ragged2e, xurl, TeX Live 2024 ou mais recente
\documentclass[12pt, a4paper, brazilian]{abntex2}
\usepackage[T1]{fontenc}
\usepackage{csquotes}
\setlrmarginsandblock{6cm}{6cm}{*}
\checkandfixthelayout
%%BEGIN:alinhamentopreambulo
\usepackage[backend=biber, style=abnt, uniquename=init, giveninits,
            repeatfields, noslsn]{biblatex}        % sem a opção justify
\usepackage{ragged2e}                             % \RaggedRight: à esquerda, com hifenização
\usepackage{xurl}                                 % URLs quebram em qualquer caractere
%%END:alinhamentopreambulo
\addbibresource{catalogo.bib}
\AtBeginBibliography{\SingleSpacing}
\hypersetup{hidelinks}                 % sem as molduras dos links (só para a comparação)

\begin{document}

\nocite{ex-artigo-online, ex-legislacao, ex-site, ex-repositorio, ex-relatorio}

\begingroup
\toggletrue{justify}                               % o mesmo que a opção justify
\printbibliography[heading=subbibliography, title={Com justify}]
\endgroup

\newpage
\printbibliography[heading=subbibliography, title={Padrão (raggedright)}]

\newpage
\begingroup
%%BEGIN:alinhamentorecomendado
\appto{\bibsetup}{\RaggedRight}    % substitui o \raggedright do estilo
\urlstyle{same}                    % URL na mesma fonte do texto
%%END:alinhamentorecomendado
\printbibliography[heading=subbibliography, title={Recomendado}]
\endgroup

\end{document}
